%!TEX root = thesis.tex
\chapter{Conclusion}
\label{chap:conclusion}

This chapter summarises the key findings of the thesis, provides direct answers to the
five research questions stated in Chapter~1, describes limitations and threats to
validity, and identifies the most important directions for future research.

\section{Summary of Findings}
\label{sec:conc:summary}

This thesis developed and evaluated methods for classifying wind turbine operational
zones from SCADA data, with the goal of enabling reliable, context-aware condition
monitoring in a PHM digital twin. The work was structured as two connected studies.

The first study (Chapter~3) defined nine operational zones based on turbine physics and
trained RIPPER and FOIL rule-based classifiers to identify these zones from four SCADA
variables. Eight of the nine zones are present in the dataset; the ninth (high-wind
safety shutdown above 25\,m/s) does not occur in the data used in this thesis. RIPPER
achieved 92.6\%
accuracy on internal validation and 91.2\% on an independent 2017 test set. FOIL achieved
92.3\% and 91.4\% respectively. Both classifiers generalised well across years and
produced interpretable rules that align with known turbine control logic.

The second study (Chapter~4) compared DRL agents (DQN, PPO, A2C) and supervised
classifiers (Random Forest, Gradient Boosting, MLP) on the same zone classification
task. The key finding was that a Random Forest trained on the four raw SCADA signals
achieves 100\% accuracy on both the 2016 validation set and the 2017 external test. Among
DRL agents, DQN with physics guardrails reached 97.7\% accuracy on the external test,
compared to 96.0\% for PPO and 93.4\% for A2C. Physics guardrails substantially improved
PPO and A2C performance and added a safety guarantee to DQN at no accuracy cost. A
deployment pipeline comparison showed that placing DQN before a LightGBM meta-model
reduces the number of readings flagged for human review by approximately 70\% compared
to building the meta-model directly on RIPPER outputs, while maintaining 98.9\% recall
of DQN's remaining errors with only 14 false alarms across 52,234 test readings.

\section{Answers to the Research Questions}
\label{sec:conc:rqa}

\subsection*{RQ1: How accurately and interpretably do physics-informed rule-based
classifiers identify operational zones, and how well do they generalise?}

RIPPER and FOIL both achieve around 91--93\% accuracy on both internal validation
and an independent external test set from a different turbine and a different year.
The rules are physically interpretable and can be verified against known turbine physics.
Errors are concentrated mainly in Zone~5 (free wheeling) and Zone~6 (generator
engagement), which share a very narrow RPM band, making them difficult to separate with
discretized bin features. Both classifiers generalise well across years.

\subsection*{RQ2: How much accuracy improvement do supervised classifiers achieve over
rule-based methods, and what explains this improvement?}

Supervised classifiers substantially outperform rule-based methods. A Random Forest
achieves 100\% accuracy on both evaluation sets, compared to 91--93\% for the
rule-based models. Gradient Boosting reaches 99.9\% and the MLP reaches 99.4\%.
The accuracy gap is explained entirely by the discretization constraint of the
rule-based methods: supervised classifiers work directly with continuous SCADA values
and can find the exact zone boundaries, while RIPPER and FOIL are limited to coarser
bin boundaries.

\subsection*{RQ3: How effectively do DRL agents learn to classify operational zones,
and under what conditions are they preferable to supervised classifiers?}

DRL agents learn this task effectively. DQN with guardrails reaches 97.7\% accuracy,
a result that is 6--7 percentage points better than the rule-based methods and only
2.3 points below the best supervised classifier. DRL is preferable in settings where
labelled training data is not available and the reward signal can be computed from
physics rules, or where the model needs to update continuously during deployment.
Where labelled data is available and accuracy is the primary concern, a supervised
classifier is the better choice.

\subsection*{RQ4: What is the contribution of RIPPER meta-features to accuracy?}

For supervised classifiers, the RIPPER meta-features contribute essentially nothing. The
Random Forest on four raw SCADA signals achieves 100\%, the same as the RF on all 13
features (which makes one error on the 2016 validation set). The MLP on raw SCADA
achieves 99.5\%, marginally better than the MLP on full features (99.4\%). The raw SCADA
signals already contain all the information needed to identify the zones correctly. The
RIPPER meta-features may still add value for DRL agents as part of what the agent observes at each step, but this was not directly tested in this study.

\subsection*{RQ5: To what extent do physics guardrails improve DRL agent accuracy,
and how does the impact vary across DRL algorithms?}

The impact of physics guardrails varies substantially across algorithms. For A2C,
guardrails improve macro-F1 by 0.269 (from 0.578 to 0.847), a large gain that makes
A2C usable where it otherwise fails. For PPO, the gain is meaningful at 0.044 (from
0.910 to 0.954). For DQN, the improvement is negligible at 0.004 (from 0.963 to 0.967)
because DQN already respects the physical constraints in 99.94\% of predictions without
guardrails. This shows guardrails are most valuable for weaker DRL agents and add a
safety guarantee to well-trained agents at essentially no cost.

\section{Limitations and Threats to Validity}
\label{sec:conc:limits}

\paragraph{Internal Validity.}
Internal validity concerns whether the comparisons between models are fair and confounding
factors have been controlled. All models are trained and evaluated on the same dataset
using the same chronological split and the same evaluation metrics, which ensures a
consistent basis for comparison. The DRL agents and supervised classifiers share the same
normalised feature set, so differences in accuracy cannot be attributed to data
preprocessing differences. One potential internal validity concern is that DRL agents are
evaluated using the best checkpoint saved during training rather than the final model
weights; this provides a conservative estimate of each algorithm's peak performance but
introduces a dependency on when during training the best checkpoint was saved.

\paragraph{External Validity.}
External validity concerns whether the findings generalise beyond the specific data and
conditions studied. The most significant limitation is that all results derive from a
single offshore wind farm with one turbine type. Whether the nine-zone framework and the
classification results transfer to onshore turbines, different manufacturers, or turbines
with fundamentally different control systems is unknown. The 100\% Random Forest accuracy
result in particular may not replicate if the zone boundaries require significant
modification for a different turbine design. The cross-turbine, cross-year external test
on Wind Turbine~11 (2017) provides evidence of generalisation within the same wind farm,
but this represents a limited form of external validity.

\paragraph{Construct Validity.}
Construct validity concerns whether the measurements used capture what the thesis claims
to measure. Overall accuracy and macro-F1 are valid measures of zone classification
correctness. However, both metrics evaluate classification as an end in itself rather
than measuring downstream PHM outcomes. A complete evaluation would test whether correct
zone classification reduces false alarm rates in fault detection, improves power curve
model accuracy, or enables better maintenance scheduling. This end-to-end evaluation is
not performed and is identified as the most important direction for future work.

\paragraph{Conclusion Validity.}
No formal statistical significance tests are applied to the accuracy differences in
Table~\ref{tab:ch4:results}. Given the external test size of 52,234 samples, a 1\%
accuracy difference corresponds to approximately 500 observations, making the larger
reported differences almost certainly significant in a statistical sense. However, a
McNemar test or bootstrap confidence interval analysis would make this claim explicit.
Small differences between models, particularly those below 0.5 percentage points, should
be interpreted cautiously without formal significance testing.

\section{Future Work}
\label{sec:conc:future}

The most important next step is multi-turbine validation. Applying the zone framework
from Chapter~3 to additional turbines beyond those used in this thesis and to turbines
from different wind farms would test how well the physics-informed definitions generalise.

Online DRL adaptation is the second priority. Deploying a DQN agent on a simulated
turbine and allowing it to update its policy from reward signals during operation would
directly test the main theoretical advantage of DRL over supervised classifiers. A reward
function that covers all zones using physics rules, not just the two clearest zones
covered by the guardrails, would be a key component of this experiment.

End-to-end PHM evaluation is the third direction. Rather than evaluating zone
classification accuracy alone, future work should build a zone-aware power curve model
and fault detection system, and measure whether correct zone identification actually
reduces false alarms and improves fault detection rates. This would directly demonstrate
the value of the zone classification framework for the PHM digital twin application that
motivated this thesis.

Finally, a targeted experiment removing the RIPPER meta-features from the DRL state
would clarify their role. Training DQN with and without the RIPPER probability features in
its observation would show whether these features help the agent learn faster or achieve
higher accuracy, or whether the four raw SCADA signals are sufficient for DRL as well as
for supervised classifiers.
